% compendium/execution-trading-costs/smart-order-routing.tex
% Entries migrated on 2026-10-15 from the legacy smart-order-routing file (slice 14), owned by
% market-structure-execution-quality; dispositions and source checks in notes/migration.md, the blocks
% and the original table verbatim in notes/artefacts/market-making-lob/derivations.tex.

\section{Smart order routing across fragmented venues}

\begin{qfdefinition}[Smart order routing: splitting a marketable order across venues, and why displayed liquidity is only probable]\label{def:smart-order-routing}
\qffield{Statement} When several venues trade the same instrument, marketability is venue-specific: with 50 shares offered at 10.00 on one venue and 150 at 11.00 on another, a buy limit at 10.50 is marketable on the first only \citep[sec.~3.4.2, pp.~233--234]{alpha-laruelle}. A smart order router (SOR) decides where the child orders of a trading algorithm go; in Lehalle and Laruelle's architecture it is one of the tactics working inside the envelope set by the risk-control layer (\hyperref[def:trading-envelope]{trading envelope}). Their example (Table~3.3, p.~235, page image) is a market buy of 500 shares against three European books, routed by three brokers (pp.~235--236):
\begin{itemize}
\item B1 has Euronext and Chi-X and splits for the best average price: 101.15;
\item B2 has all three venues but sends the whole order to the single venue with the best price: the text says 102.50 on Chi-X;
\item B3 has all three and splits across them: 100.75, with a free choice of venue, ``surely'' by fees, for the last 50 shares at 102.00.
\end{itemize}
With latency or duplicated (``mirrored'') liquidity the ranking changes. If a high-frequency trader co-located at Euronext and Chi-X mirrors 150 shares on Turquoise, which the brokers reach 2~ms later than Euronext, it cancels them before B3's order arrives; B3 learns of the 150 unfilled shares about 6~ms later and, with prices rising, ends up worse than B1. What matters is not the visible liquidity but the probability of obtaining it at a price when the order arrives, given latency differences, mirrored orders and hidden liquidity (pp.~238--239).

\qffield{Notation} Interpretation, for the routing decision at the time a child order is released: the inputs are the visible books, the venues the broker's policy allows, the latencies to each, fees, and estimates of hidden liquidity and fill rates; the decision is a quantity per venue (and, for passive orders, a price); the constraints are the child order's size, side and limit and the parent's envelope. The SOR does not decide how much the parent trades when; that is the schedule (\hyperref[mod:almgren-chriss-execution]{mean--variance execution}).

\qffield{Conditions} The example assumes no latency issue for the base case (p.~235) and full access to each venue's displayed depth. Liquidity seeking for a marketable order must allow for hidden liquidity and for the chance that a fast trader cancels when it sees the order reach another venue first; limit orders raise further issues (p.~239). A liquidity seeker removes liquidity using estimates of hidden liquidity (including dark pools), mirrored orders and latency, and posts limit orders to finish fast; it is an ``enhanced limit order'' lasting from about 10 minutes in liquid stocks to a few hours in illiquid ones, while adverse selection from finishing fast is left to a scheduling layer at 10 minutes to a few hours; the hard part is estimating fill rates and real liquidity, which cluster (pp.~242--243).

\qffield{Practice} Lehalle and Laruelle allocate across venues from ``fill rate curves'', the expected fill intensity against distance from the best opposite price, which differ by venue (Figure~3.17, pp.~243--244), citing Laruelle et al.\ (2011, 2013), not consulted. Passive splits are in \hyperref[meth:passive-order-split]{passive order split}; fee effects in \hyperref[res:fee-adjusted-venue-choice]{fee-adjusted venue choice}.

\qffield{Pitfalls} Interpretation, recomputed from Table~3.3 by sweeping the cheapest levels first: B1 101.15 and B3 100.75 are right; the best single venue for all 500 shares is Chi-X at 102.05 (Euronext 102.10, Turquoise 102.20), so B2's ``102.50'' looks like transposed digits. The bracketed quantities for B1, $(25+50+100+100)$ and $(0+25+100+200)$, add to 600 and list depth up to 102.00, not what was bought (200 of the 300 shares at 102.00). The legacy sweep taxonomy's ``local best price: cheapest visible venue first'' is not B2's rule, which sends the order to one venue. Its routing-logic list (queue position, reject history, gateway and matching-engine delays, information cost of probing) and its four-step routing loop that updates venue scores from fills are a reasonable synthesis but not lists in the source; only fees, latency, hidden and mirrored liquidity and fill-rate estimates are. Displayed depth cannot be assumed to stay available while a sequence of child orders is routed. Venue names and the HFT example describe European markets around 2018.

\qffield{See also} \hyperref[def:trading-envelope]{trading envelope} (routers as tactics; owned by \texttt{execution-trading-costs}); \hyperref[def:order-exposure-decisions]{order exposure decisions} (whether to show and where); \hyperref[def:best-execution]{best execution} (routing duties and price protection); \hyperref[meth:passive-order-split]{passive order split}; \hyperref[res:fee-adjusted-venue-choice]{fee-adjusted venue choice}.

\qffield{Further reading} The legacy blocks it replaces, verbatim: \nolinkurl{notes/artefacts/market-making-lob/derivations.tex}.

\qffield{Index} smart order routing; SOR; marketable price; marketable order; fragmentation; venue choice; child order; sweep; latency; mirrored liquidity; duplicated quotes; hidden liquidity; fill rate curves; liquidity seeker; liquidity seeking; co-location

\qffield{Status} Verification: partly checked 2026-10-15 (against the text layer of a local copy of \texttt{alpha-laruelle}, 2018 World Scientific printing, sec.~3.4.2, pp.~233--244, page images of pp.~235 and~236; the example recomputed). Stability: draft. Currency: European venues and latencies as of the 2018 printing.
\end{qfdefinition}

\begin{qfresult}[Fee-adjusted venue choice: a worse price with lower fees wins only when the price step is smaller than the fee difference]\label{res:fee-adjusted-venue-choice}
\qffield{Statement} For a buy of $Q$ shares, taking liquidity on venue $V$ at price $P$ with proportional fee $f$, against venue $V'$ at $P(1+\delta)$ with fee $f-\phi$, $V'$ is cheaper when
\[ P(1+\delta)(1+f-\phi)Q<P(1+f)Q\iff\delta<\Bigl(1-\frac{\phi}{1+f}\Bigr)^{-1}-1\simeq\phi, \]
so a worse price is worth taking only if the tick, in relative terms, is about as small as the fee difference: at a price of 100 euros and a fee difference of 0.1 basis point, the tick would have to be below 0.1 cent \citep[sec.~3.4.2, pp.~236--237, page image of p.~236]{alpha-laruelle}. Figure~3.16 plots the critical tick against the fee for $\rho=\phi/f$ (p.~237). Liquidity providers, who pay a fee or receive a rebate when their limit orders fill, would post first on the venue cheapest for them, and market orders would go first to the venue cheapest for takers; if these differ, both cannot persist in equilibrium (pp.~237--238).

\qffield{Notation} $f$ and $\phi$ are fees as fractions of traded value, paid by the taker; a rebate is a negative fee. $\delta$ is the relative price difference, at least one tick.

\qffield{Conditions} Interpretation: the threshold is exactly $\delta<\phi/(1+f-\phi)$ for a buy, comparing acquisition cost. For a sale, comparing proceeds, $V'$ at $P(1-\delta)$ with fee $f-\phi$ is better when $P(1-\delta)(1-f+\phi)>P(1-f)$, that is $\delta<\phi/(1-f+\phi)$, also about $\phi$. Illustrative values $P=100$, $f=3$~bp, $\phi=1$~bp: the buy threshold is $0.99980$~bp and the sell threshold $1.00020$~bp; at $\delta=0.99$~bp $V'$ wins on both sides, at $1.01$~bp $V$ wins (checked with exact arithmetic). With fees per share, $c$ on $V$ and $c'$ on $V'$, the condition is $P\delta<c-c'$ on both sides: a one-cent tick beats a fee saving of 0.2 cents. The comparison assumes both quotes are executable for the whole quantity at the time of routing, both orders take liquidity, and the trader actually pays the venue fees; the spread is inside the two prices and is not added again.

\qffield{Practice} Interpretation: whether a broker passes venue fees and rebates to the client decides whose cost this is; when it does not, the comparison drives the broker's net, not the client's (\hyperref[def:broker-agency-conflicts]{broker conflicts}). A favourable fee-adjusted quote says nothing about the chance of filling there, or about later price moves.

\qffield{Pitfalls} The legacy caveat stated the inequality correctly but called it ``price priority''; it is a cost comparison between two takers' prices, not a priority rule. Current fee levels and caps are not checked here; access-fee caps belong to this pathway's reading \texttt{ex-sec-tick-access-2024}, not read.

\qffield{See also} \hyperref[def:smart-order-routing]{smart order routing}; \hyperref[def:transaction-costs]{transaction costs} (fees as explicit costs); \hyperref[def:broker-agency-conflicts]{broker conflicts}.

\qffield{Further reading} The legacy block it replaces, verbatim: \nolinkurl{notes/artefacts/market-making-lob/derivations.tex}.

\qffield{Index} fee-adjusted price; trading fees; maker--taker; taker fee; rebate; tick size; critical tick; venue choice; fee arbitrage; access fees

\qffield{Status} Verification: partly checked 2026-10-15 (against the text layer of a local copy of \texttt{alpha-laruelle}, 2018 World Scientific printing, sec.~3.4.2, pp.~236--238, page image of p.~236; buy and sell thresholds derived and checked). Stability: draft. Currency: timeless (arithmetic); no current fee schedule checked.
\end{qfresult}

\begin{qfmethod}[Passive order split: ranking venues by queue, distance and trading rate, and the minimum-resting-time split]\label{meth:passive-order-split}
\qffield{Statement} To post a buy limit of $D$ shares below the best ask across venues, Lehalle and Laruelle compare three criteria on a stylised book (Table~3.4, p.~240, page image): rank in the queue (shares ahead at the limit price), distance of the venue's best ask from the limit, which shows where sellers are aggressive, and the recent selling rate at that price. With $Q_n$ shares ahead and selling rate $r_n$ on venue $n$, an order of size $D$ takes on average
\[ \tau_n(D)=\frac{D+Q_n}{r_n} \]
to fill (eq.~(3.3)). Splitting $D$ so that every part finishes at the same time $\tau^\star$ gives $f_nD=\tau^\star r_n-Q_n$ and, summing, $\tau^\star=(D+\sum_nQ_n)/\sum_nr_n$ \citep[sec.~3.4.2, pp.~239--241, n.~16, page images of pp.~240--241]{alpha-laruelle}. In their example ($D=500$; $Q=450, 100, 200$; $r=10, 5, 5$ shares a second on Euronext, Chi-X and Turquoise), $\tau(500)$ is 95, 120 and 140~s, $\tau^\star=62.5$~s, and the split is 35, 42.5 and 22.5 percent, 175, 212 and 113 shares (Table~3.5, p.~241). The criteria rank the venues differently (Table~3.6, p.~242):

\smallskip
\noindent{\small\begin{tabularx}{\linewidth}{@{}>{\RaggedRight}p{0.3\linewidth}>{\RaggedRight}X>{\RaggedRight}X>{\RaggedRight\arraybackslash}X@{}}
\toprule
Criterion & 1st venue & 2nd venue & 3rd venue\\
\midrule
Rank in the queue & Chi-X & Turquoise & Euronext\\
Distance to the best & Euronext & Turquoise & Chi-X\\
Trading rate & Euronext & Chi-X or Turquoise (tied) & \\
Minimum resting time (split) & Chi-X (212) & Euronext (175) & Turquoise (113)\\
\bottomrule
\end{tabularx}}
\smallskip

\noindent Sources: rows from \citep[sec.~3.4.2, Table~3.6, p.~242]{alpha-laruelle}, with the split sizes from Table~3.5; checked by recomputation.

\qffield{Conditions} Interpretation: the formula treats the queue ahead and the selling rate as fixed and known, ignores cancellations ahead of the order and later arrivals, and fills each part at the average rate, so $\tau$ is a deterministic approximation; the source points to stochastic methods for accurate splits (Laruelle et al.\ 2011, not consulted). Equal finishing times minimise the time to complete the whole order, the latest part, rather than an average.

\qffield{Practice} Interpretation: a venue with $f_n<0$ cannot finish by $\tau^\star$ even with nothing posted; drop it and recompute $\tau^\star$ over the rest. Setting $f_n$ to zero without recomputing leaves fractions that no longer sum to one. Lehalle and Laruelle's liquidity seeker combines this posting with optimal removal of liquidity and keeps the parent inside a scheduling layer (p.~242; \hyperref[def:smart-order-routing]{smart order routing}).

\qffield{Pitfalls} The book prints eq.~(3.4) as $\tau^\star=(D-\sum_nQ_n)/\sum_nr_n$ (p.~241, page image), but its own footnote gives $D=\tau^\star\sum_nr_n-\sum_nQ_n$ and Table~3.5's 62.5~s needs $D+\sum_nQ_n$; with the minus sign the example gives $-12.5$~s. The legacy definition has the plus sign. The text calls the Euronext rate ``buy orders hit the 95.00 price level''; for a resting buy order the relevant flow is sell orders, as for the other venues. The legacy table generalised Table~3.6 without venue names (``highest adjusted rate'', ``residual venue''), which does not match the source's rows, and the legacy taxonomy's ``resting-time target'' is the split, not a separate criterion. The split sizes are rounded (212.5 and 112.5).

\qffield{See also} \hyperref[def:smart-order-routing]{smart order routing}; \hyperref[def:order-exposure-decisions]{order exposure decisions}; queue priority and limit-order fill risk in the microstructure overview's entries.

\qffield{Further reading} The legacy blocks and table, verbatim: \nolinkurl{notes/artefacts/market-making-lob/derivations.tex}.

\qffield{Index} passive order split; limit order placement; queue position; rank in queue; distance to best; trading rate; minimum resting time; fill time; venue ranking; Lehalle--Laruelle Table 3.6

\qffield{Status} Verification: partly checked 2026-10-15 (against the text layer of a local copy of \texttt{alpha-laruelle}, 2018 World Scientific printing, sec.~3.4.2, pp.~239--242, page images of pp.~240 and~241; the example and the corrected formula recomputed with exact arithmetic). Stability: draft. Currency: timeless (method); the example is stylised.
\end{qfmethod}
